% !TEX root = ../main.tex
% =====================================================================
%  §10  物理进路的诸应用：后测
%  §11  物理进路的诸应用：预测
%  §12  可能的未来方向
%  原文 pp.67--71（PDF pp.68--72）
%  公式：(10.1)、(11.1)--(11.8)，逐条保留原文编号
%  图占位 4 处：原 PDF 中均只有一行图题、并无实际图形
% =====================================================================

\section{物理进路的诸应用：后测}

\begin{keytakeaways}
\item \textbf{物理图像：} 塞伯格--威滕方程并非一套孤立的新方程，而是超导的朗道--金兹堡低能有效理论在自对偶、临界耦合处的孪生体；自旋场的零点集扮演磁通线（涡旋）核心的角色，其“世界面”就是伪全纯曲线。于是辛流形上的 SW 不变量与计数伪全纯曲线的格罗莫夫--威滕不变量必然重合。
\item \textbf{数学工具：} 伪全纯曲线与 GW 不变量、芬图谢尔--斯特恩 / 格策--扎吉尔（格策--扎吉尔）爆破公式、$u$ 平面上的算子乘积展开 $\exp[t\mathcal{O}(E)]=\sum_{k\ge1}t^{k}B_{k}(\mathcal{O})$，以及非单连通流形上的赖德迈斯特--米尔诺挠。
\item \textbf{核心结论：} 物理进路把若干已知数学定理重新解释为“后测”：爆破公式成为例外曲面上 2-可观测量的算子乘积展开，辛流形上的 $\mathrm{SW}=\mathrm{GW}$ 成为磁通线世界面的计数；但如作者所强调，数学家迄今并未真正动用这些物理洞见。
\end{keytakeaways}

塞伯格--威滕（塞伯格--威滕）方程被引入四维流形的理论，带来了直接而巨大的冲击，并且至今仍在深刻地影响这一领域。此前数学家用了多年时间建造起一整套技术武库，以处理比它困难得多的非阿贝尔方程；而在 SW 方程问世之后，唐纳森（Donaldson）的所有定理都在数周之内被重新证明。不仅如此，SW 方程还导致若干长期悬而未决的猜想（例如托姆（Thom）猜想）获得解决，并使数学家的视野超出了唐纳森理论已达到的范围。

% [Tutor Note: “后测”（postdiction）是作者自造的词，指先有数学定理、事后才由物理图像给出解释的情形——与“预测”方向相反。读本节须始终留意这个时间箭头：作者并不是用物理预言了陶布斯定理，而是事后指出它本来可以被物理预见。这也是 §10.1 末尾那句“对物理学家而言不幸的是，这只能算一次物理上的后测”的由来：若历史顺序颠倒，SW 不变量就会像 GW 不变量一样是被“算出来”的。]

欲了解清晰而精到的叙述，可参看 S. Donaldson（Donaldson）的综述 \cite{ref9}；关于这段历史的通俗介绍见 \cite{ref19,ref20}。

总的来说，数学家既已看出可以用阿贝尔规范理论中技术上易于处理的 SW 方程取代非阿贝尔的反自对偶方程，实际上却并没有真正使用我方才解释的那些物理洞见。尽管如此，还是出现了若干数学上的“后测”——即对已知数学事实所作的有启发性的物理解释；也出现了若干“预测”——由物理动机驱动的新数学命题，它们定义得足够明确，因而理应能够承受严格的数学证明（或反例）。

下面我们列出若干数学后测，并在下一节简要回顾一些预测。

\subsection{辛流形上的 $\mathrm{SW}=\mathrm{GW}$}

在一些漂亮的工作中，C. Taubes（Taubes）证明：\emph{辛流形}上的 SW 不变量可以等同于计数伪全纯曲线的格罗莫夫--威滕（格罗莫夫--威滕）不变量。这一点本来可以从超导物理中预见到，因为 SW 方程与超导的朗道--金兹堡（Landau--Ginzburg）低能有效理论的方程非常相似。这里所说的伪全纯曲线，在物理上可以看作阿布里科索夫--戈尔科夫（Abrikosov--Gorkov）磁通线的世界面。关于这一点，可参看威滕（Witten）的漂亮文章 \cite{ref57}。对物理学家而言不幸的是，这只能算一次物理上的后测，尽管历史本来很容易以另一种顺序展开。

% [Tutor Note: 此处“非常相似”其实是逐项对应，值得点明：SW 方程为
%  $F_{A}^{+}=\tfrac12(\bar\psi\psi)_{0}$ 与 $D_{A}\psi=0$（下标 $0$ 表示 2-形式的无迹部分）。
%  把自旋场 $\psi$ 换成复标量序参量（即 GL 的希格斯场）、把 $U(1)$ 规范场换成超导电流的
%  规范场，第一条就变成 GL 的超导电流关系，第二条则是把 GL 的二阶方程降阶所得的一阶
%  （Bogomolnyi）方程。GL 参数取临界值 $\kappa=1/\sqrt2$ 恰是使 GL 泛函可配成平方和的取值；
%  这与 SW 泛函可由 Weitzenböck 型配平方得到完全平行。换言之，二者共享同一个“BPS 平方和”
%  结构，而不是仅在形式上相像。]

% [Tutor Note: 磁通线 $\leftrightarrow$ 伪全纯曲线的机制可以这样读：二维 GL 的涡旋解中，
%  序参量在某一点为零、绕该点一周的相位绕数为整数 $n$，磁通量子化 $\int F=2\pi n$；
%  一阶涡旋方程在复坐标下化为 $\bar\partial$ 型方程，故“涡旋核心”满足 Cauchy--Riemann
%  方程。把二维涡旋沿时间演化，核心扫出的世界面在四维中就是伪全纯曲线。陶布斯的定理
%  把这一图像精确化：把辛形式的大参数缩放加进 SW 方程后，解集中于伪全纯曲线附近，
%  $F^{+}$ 的流即该曲线的面积形式。还须注意，这条对应只对 $\lambda=\pm\lambda_{c}$ 这类
%  规范类（即 $\lambda^{2}=2\chi+3\sigma$ 的类）成立：辛流形上的 SW 基本类只有 $\pm\lambda_{c}$。]

\subsection{爆破公式}

这一方法给出了芬图谢尔--斯特恩 / 格策--扎吉尔（格策--扎吉尔）爆破公式的一个简单物理推导：该公式可看作爆破的例外曲面上 2-可观测量的某种“算子乘积展开”。从拓扑上说，一次爆破就是 $\hat X=X\#\CC\PP^{2}$，而由 $u$ 平面容易看出
\begin{equation}
  \exp\bigl[t\mathcal{O}(E)\bigr]=\sum_{k=1}^{\infty}t^{k}B_{k}(\mathcal{O}),
  \tag{10.1}
\end{equation}
其中 $B_{k}(\mathcal{O})$ 是多项式。这是二维共形场论中那个熟悉做法的四维版本：用一个局部算子的无穷和来替换一个把手。细节见 \cite{ref38}；更高阶情形的推广见 \cite{ref10,ref24}。

\begin{figure}[H]
  \centering
  \figplaceholder{〔原文此处于第 68 页（PDF 第 69 页）留有图位；原 PDF 中未插入实际图形〕}
  \caption{与“在一个把手上作因子化”的比较。}
  \label{fig:handle-factorization}
\end{figure}
% [原 PDF 第 69 页]
% [Tutor Note: 原 PDF 第 69 页此处只有一行大写图题 “FIGURE: COMPARISON WITH FACTORIZATION
%  ON A HANDLE”，并没有插入任何图形；本译本按规范以占位框重现该图位，绝不臆造图形。
%  图题本身也只是提示性的，没有给出坐标轴或标注。]

% [Tutor Note: 称 (10.1) 为“算子乘积展开”，其准确含义是：连通和 $X\#\CC\PP^{2}$ 的配分函数
%  并不等于两个因子配分函数之积，而是在 $X$ 的 $u$ 平面积分上插入一个只与例外曲面
%  $E$（$E^{2}=-1$）有关的 2-可观测量 $\mathcal{O}(E)$ 的指数；把指数展开即得 (10.1)，
%  其中每一项 $t^{k}B_{k}(\mathcal{O})$ 都是 $\mathcal{O}$ 的多项式，即 $p$、$\Sigma$ 的多项式
%  （等价地，是自旋$^{c}$ 类 $\lambda$ 的多项式）。因此“因子化”在四维只能成立到
%  “乘以一个算子插入”的程度。]

% [Tutor Note: 二维类比值得展开：在二维共形场论中，往带边 Riemann 面上粘一个把手，等价于在
%  原曲面上插入一个态，而后者可按局部算子的无穷和展开，“带把手”的关联函数等于
%  $\sum_{k}t^{k}\langle\cdots\mathcal{O}_{k}\rangle$（$\langle\cdots\rangle$ 表示其余插入），
%  其中 $t$ 是把手的模参数。四维中没有真正的局域算子代数，但 $u$ 平面积分只用到
%  2-可观测量的关联函数，而“与 $\CC\PP^{2}$ 作连通和”这一操作恰好可用
%  $\sum_{k\ge1}t^{k}B_{k}(\mathcal{O})$ 代替，故 (10.1) 正是该二维手法的四维对应物。
%  另需注意求和自 $k=1$ 起而非 $k=0$：$k=0$ 项只是常数重标度，原文未说明它的去向，
%  按 [38] 的处理它应已并入未写出的整体归一化因子。]

\subsection{孟--陶布斯：SW 不变量与赖德迈斯特--米尔诺挠}

$Z_{DW}$ 的推导可以推广到非单连通的四维流形；就纯 $SU(2)$ 矢量多重态理论而言，这已在 \cite{ref30,ref32} 中完成。\footnote{〔原稿边注〕此处应补充说明该文的内容。}

% [Tutor Note: 非单连通情形有三处须小心。(i) 自旋$^{c}$ 结构此时构成 $H^{2}(X;\ZZ)$ 的一个
%  torsor，而不再只是格 $\Gamma_{w}$ 的平移，故 (9.1) 一类的求和要换成对整个 torsor 求和。
%  (ii) $b_{1}>0$ 时库仑分支不再是单纯的 $u$ 平面，而应再乘上环面
%  $H^{1}(X;\RR)/H^{1}(X;\ZZ)$，$u$ 平面积分因此多出零模（通量扇区）上的积分。
%  (iii) 希格斯分支的模空间会出现可约解。标题所引的孟--陶布斯定理则属于三维情形：
%  三维流形的 SW 不变量等于该流形的赖德迈斯特--米尔诺（等价地 Turaev）挠。]

\section{物理进路的诸应用：预测}

\begin{keytakeaways}
\item \textbf{物理图像：} 若把 $SU(2)$ SYM 耦合到 $N_{f}\le4$ 个超多重态并赋予复质量参数，则 $u$ 平面上的塞伯格--威滕曲线族会在 $2+N_{f}$ 个点退化；当这些退化点随质量参数移动而相互碰撞时，理论走到超共形点，$u$ 平面积分可能出现发散。
\item \textbf{数学工具：} 虚二次域类数与模拟模形式、$\Theta$ 提升、$SU(N)$ 的高阶唐纳森不变量、地理问题的 $(\chi,\sigma,t)$ 三元组与 $c$--$\chi_{h}$ 平面、生成函数 $\mathrm{SW}_{X}(z)$ 在 $z=0$ 的零点阶判据。
\item \textbf{核心结论：} 要求 $Z_{DW}$ 在超共形碰撞极限下有限，即“超共形单型”判据，由此得到基本类数目的下界 (11.6) 与新的地理界线——广义诺特界 $c\ge\chi_{h}-2B-1$ (11.7)。这些是真正意义上的预测：它们是可以被数学证明或否证的新命题。
\end{keytakeaways}

最引人注目的物理预测正是我们前面反复强调的那一条：对 $b_{2}^{+}>1$ 的流形，$Z_{DW}$ 可以只用 SW 不变量写出。

不过仍值得一问：对四维流形的数学，是否还有别的物理预测？这里我们简要提及几项。\footnote{〔原稿边注〕去查孟--Chwan--Tan 的几篇文章；他究竟做了什么？}

\subsection{类数的新公式}

在若干“好”的情形下，$u$ 平面积分确实可以显式算出。最值得注意的是 $\CC\PP^{2}$ 的答案，它揭示出与虚二次域的类数以及模拟模形式（mock modular form）之间的一种耐人寻味的关系。事实上，$u$ 平面积分与某些类型的“$\Theta$ 提升”密切相关，而后者在数论与弦微扰论中都曾出现；在后一个语境中，它们被用来给出 Borcherds 关于自守积（automorphic product）的结果的概念性证明。\footnote{〔原稿边注〕此处原稿只写了占位标记 “[Cite: HarveyMoore]”，未给出具体文献编号。}

% [Tutor Note: “模拟模形式”的出现是结构性的，并非巧合：$\CC\PP^{2}$ 的 $u$ 平面积分并不是对
%  整个模基本域的积分，而是“截断”在一个带边区域上；边界（例如尖点 $u=\infty$ 处）贡献一个
%  非全纯的修正项，使结果只满足一个带“阴影”的模变换律——这正是扎吉尔意义下模拟模形式的
%  定义。虚二次域类数 $h(-n)$ 的生成函数是模拟模形式的经典范例，故这里的数论关系可以
%  理解为：$u$ 平面积分自动把计数问题组织成类数的生成函数。]

\subsection{其他单规范群的唐纳森不变量}

如上面所强调的，前述技术可以相当直接地推广到高阶不变量。对 $SU(N)$ 的唐纳森不变量，存在上述公式的类比 \cite{ref31}。令人失望的是，它再一次完全可以由经典上同调环与塞伯格--威滕不变量表出。克龙海默（Kronheimer）已对某些特殊的 $X$ 验证了这一预测 \cite{ref22}。\footnote{〔原稿边注〕此处应大幅补充：给出一般 $G$ 的威滕猜想的类比。}

四维流形的拓扑与超共形不动点的存在性之间的联系，导致了拓扑学中一些非平凡的新结果 \cite{ref33,ref34}。

% [Tutor Note: “令人失望”恰恰是信息所在。高阶唐纳森不变量的定义涉及秩 $N-1$ 的库仑分支、
%  多重覆盖以及穿墙现象，若这样复杂的对象最终仍只依赖经典上同调环与 SW 不变量，则它不提供
%  新的微分拓扑信息。从物理上看这可以预期：拓扑扭曲的 $\mathcal{N}=2$ 理论的红外数据已被
%  SW 不变量完全编码。反过来说，克龙海默的验证（对某些 $X$）也说明：高阶公式的符号与
%  归一化必须与高阶情形的穿墙公式相容，这是一项非平凡的检验。]

\subsection{地理问题}

对一个紧致四维流形，我们可以联系三元组 $(\chi,\sigma,t)\in\ZZ\times\ZZ\times\ZZ_{2}$，其中 $t$ 是“型”，告诉我们相交形式是偶的还是奇的。

地理问题问的是：哪些取值可以实现？以及对给定的 $(\chi,\sigma,t)$，有多少个例子（即互不同胚的四维流形）？出色的综述见 \cite{ref50,ref12}。

关于唯一性，显然需要加一些限制以避免平凡情形。例如
\begin{equation}
  n\CC\PP^{2}\#\, m\overline{\CC\PP}^{2}
  \tag{11.1}
\end{equation}
有 $\chi=2+m+n$ 与 $\sigma=n-m$；又因 $\chi+\sigma=2(1+b_{2}^{+}-b_{1})$ 是偶数，除 $m\ge0$、$n\ge0$ 之外，对 $\chi,\sigma$ 本质上没有任何限制。

于是可以考察：
\begin{enumerate}[label=\arabic*., leftmargin=2.1em]
\item 复流形；
\item 辛流形；
\item 不可约流形（意思是：$X=X_{1}\#X_{2}$ 蕴含 $X_{1}$ 或 $X_{2}$ 是 $S^{4}$）。
\end{enumerate}

把这些界用如下变量画出最为方便：
\begin{equation}
  c:=2\chi+3\sigma,\qquad
  \chi_{h}:=\frac{\chi+\sigma}{4}=\frac{1+b_{2}^{+}-b_{1}}{2}.
  \tag{11.2}
\end{equation}
$\chi_{h}$ 可以是整数或半整数。若它是整数，则 $X$ 容许一个近复结构；此时 $c=c_{1}(X)^{2}$。若 $X$ 是复流形，则 $\chi_{h}$ 就是全纯欧拉示性数。

% [Tutor Note: (11.1) 中的计数值得逐项核对：$\chi(\CC\PP^{2})=3$、$\sigma(\CC\PP^{2})=\pm1$，
%  连通和满足 $\chi(X\#Y)=\chi(X)+\chi(Y)-2$、$\sigma(X\#Y)=\sigma(X)+\sigma(Y)$。于是
%  $\chi=3n+3m-2(n+m-1)=2+m+n$，$\sigma=n-m$。(11.2) 中
%  $\chi_{h}=(\chi+\sigma)/4=(1+b_{2}^{+}-b_{1})/2$ 之所以成立，是因为
%  $\chi=2-2b_{1}+b_{2}^{+}+b_{2}^{-}$、$\sigma=b_{2}^{+}-b_{2}^{-}$，相加即得
%  $\chi+\sigma=2-2b_{1}+2b_{2}^{+}$。注意 $\chi$ 与 $c$ 都\emph{不}在连通和下相加：
%  一次爆破使 $(\chi,\sigma)\mapsto(\chi+1,\sigma-1)$，故 $\chi_{h}$ 在爆破下不变，而
%  $c\mapsto c-1$——这正是 $\chi_{h}$ 与 $c$ 这一对变量在爆破公式 (§10.2) 中自然出现的原因。]

% [Tutor Note: 原文“若 $\chi_{h}$ 为整数则 $X$ 容许近复结构”一句只能按\emph{必要条件}理解，
%  反向不成立。必要性的来源是整性：由阿蒂亚--辛格指标定理，近复流形的 Todd 亏格
%  $T=(c_{1}^{2}+c_{2})/12=(2\chi+3\sigma)/12=(\chi+\sigma)/4=\chi_{h}$ 必为整数。
%  反例：$X=\CC\PP^{2}\#\CC\PP^{2}$ 有 $\chi=4$、$\sigma=0$，故 $\chi_{h}=1$ 为整数；
%  但 Wu 公式要求 $c_{1}\equiv w_{2}\equiv x_{1}+x_{2}\pmod 2$，即
%  $c_{1}=ax_{1}+bx_{2}$ 且 $a,b$ 为奇数，于是 $c_{1}^{2}=a^{2}+b^{2}\equiv2\pmod 8$；
%  而近复结构要求 $c_{1}^{2}=2\chi+3\sigma=8\equiv0\pmod 8$，矛盾。所以该流形并不容许
%  近复结构。]

\begin{figure}[H]
  \centering
  \figplaceholder{〔原文此处于第 69 页（PDF 第 70 页）留有图位；原 PDF 中未插入实际图形〕}
  \caption{已知的某些界。画出一般型极小曲面：$2\chi_{h}-6\le c\le 9\chi_{h}$。}
  \label{fig:known-bounds}
\end{figure}
% [原 PDF 第 70 页]
% [Tutor Note: 这是原文唯一带图号的图题（“FIGURE1”），但原 PDF 中同样只有这一行图题、
%  并没有任何图形，故本译本只以占位框重现，绝不伪造曲线或坐标轴。图题给出的两条界即
%  下界 $2\chi_{h}-6$（诺特界 (11.8)，适用于一般型极小曲面）与上界 $9\chi_{h}$
%  （BMY 不等式，即 Bogomolov--Miyaoka--Yau 不等式）；图题中点明“一般型极小曲面”，
%  正是提醒读者这两条界只在代数曲面的这一区域内成立。]

\subsubsection{超共形奇点}

现在可以看到，超共形点的物理实际上对地理问题有所作用。

其机制如下。如我们所强调的，威滕公式有一个自然的推广：$SU(2)$ SYM 耦合到 $N_{f}$ 个超多重态，$R$ 为基本表示。UV 量子场论只在 $N_{f}\le4$ 时才是良定义的，故我们限制在这一情形。每个超多重态带一个复的“质量参数” $m_{i}$。（在数学上，$m_{i}$ 是等变上同调中的参数，这是拉巴斯蒂达（Labastida）与马里诺（Mariño）的一个结果。）\footnote{〔原稿边注〕此处原稿只写了占位标记 “[cite]”，未给出文献编号。}

量子真空模空间依旧是复平面，以 $u\in\CC$ 为参数；但此时 $u$ 平面上塞伯格--威滕族中的曲线 $\Sigma_{u}$ 会在 $2+N_{f}$ 个点 $u_{j}$（$j=1,\dots,2+N_{f}$）处退化。在每一个这样的点上，都有一种不同的 BPS 态变为无质量：

\begin{figure}[H]
  \centering
  \figplaceholder{〔原文此处于第 69 页（PDF 第 70 页）留有图位；原 PDF 中未插入实际图形〕}
  \caption{带有若干奇点的 $u$ 平面。}
  \label{fig:uplane-singularities}
\end{figure}
% [原 PDF 第 70 页]
% [Tutor Note: 原 PDF 在此处同样只有一行图题 “FIGURE OF U-PLANE WITH SEVERAL SINGULARITIES”，
%  并无实际图形；图上本应标出 $2+N_{f}$ 个奇点 $u_{j}$ 及其对应的无质量 BPS 态类型。]

对 $b_{2}^{+}>1$ 且为 SW 单型的 $X$，配分函数成为
\begin{equation}
  Z_{DW}(p,\Sigma;m_{i})
  =\tilde\alpha^{\chi}\tilde\beta^{\sigma}
  \sum_{j=1,\dots,2+N_{f}}\kappa_{j}^{\chi_{h}}
  \Bigl(\frac{\dd u}{\dd a}\Bigr)_{j}^{\chi_{h}+\sigma}
  \sum_{\lambda}\mathrm{SW}(\lambda)\,
  \ee^{\,p u_{j}+\Sigma^{2}T_{j}-\ii\bigl(\frac{\dd u}{\dd a}\bigr)_{j}\Sigma\cdot\lambda}.
  \tag{11.3}
\end{equation}
这里
\begin{enumerate}[label=\arabic*., leftmargin=2.1em]
\item $\tilde\alpha$ 与 $\tilde\beta$ 是与前面略有不同的数值常数。我们已取 $\Lambda_{N_{f}}=1$。
\item $\kappa_{j}$ 由 $u=u_{j}+\kappa_{j}q_{j}+\cdots$ 定义，其中 $q_{j}=\ee^{2\pi\ii\tau_{j}}$ 是奇点 $u_{j}$ 附近相应的模参数。
\end{enumerate}

现在，对一般的质量参数取值，$Z_{DW}(p,\Sigma;m_{i})$ 是一个明显有限且良定义的表达式。然而，当我们改变质量参数时，判别式轨迹上的点 $u_{j}$ 会移动，它们甚至可能相互碰撞。当这种碰撞涉及既带磁荷又带电荷的无质量粒子时\footnote{严格说来，是彼此非局域的（non-mutually-local）。}，会出现进一步的奇点。在数学上，这在小平（Kodaira）关于椭圆纤维化的分类中是熟悉的。

下面集中讨论 $N_{f}=1$。对 $N_{f}=1$，存在一个点 $m_{*}$，在那里两个奇点碰撞到同一个点 $u_{*}$。若令 $m=m_{*}+z$、$u=u_{*}+z+\delta u$，则塞伯格--威滕曲线在首阶为
\begin{equation}
  y^{2}=x^{3}+zx+\delta u,
  \tag{11.4}
\end{equation}
（差一些数值系数）。于是 $\kappa_{j}$ 与 $\bigl(\frac{\dd u}{\dd a}\bigr)_{j}$ 中会出现额外的零点。由于 $\chi+\sigma$ 完全可能是负的，当 $m\to m_{*}$ 时 $Z_{DW}(p,\Sigma;m)$ 可能出现发散。

然而，从物理观点看，当 $X$ 是紧致流形时不可能存在这种发散。

原因在于：在红外，奇点只能来自时空或模空间中的非紧区域。但 $X$ 是紧致的，而对 $N_{f}=1$ 不存在这样的非紧区域。（对 $N_{f}>1$，超共形奇点有时会涉及非紧区域。）

要求当 $m\to m_{*}$ 时 $Z_{DW}(p,\Sigma;m)$ 保持有限，结果会给出关于 SW 不变量的非平凡事实。

% [Tutor Note: 原文这一句 “Requiring that $Z_{DW}(p,\Sigma;m)$ as $m\to m_{*}$ turns out to imply
%  nontrivial facts about the SW invariants” 缺了谓语（应当是 “remain finite”），是一条未写完的
%  句子；本译本按上下文补足“保持有限”，特此说明。]

% [Tutor Note: (11.3) 中的 $T_{j}$ 原文未加定义。按 §7 的记号，$\Sigma^{2}T(u)$ 是接触项，
%  故 $T_{j}$ 应当就是 $T(u_{j})$；$\tilde\alpha$、$\tilde\beta$ 也只在下面的列表里被含糊地
%  称为“与前面略有不同的数值常数”。引用 (11.3) 作具体计算时，须回查 \cite{ref33} 以确认
%  归一化与 $T(u)$ 的常数项。]

% [Tutor Note: 奇点碰撞为何会在 $\kappa_{j}$ 与 $(\dd u/\dd a)_{j}$ 中制造额外零点：$\kappa_{j}$
%  是 $u-u_{j}$ 关于局部模参数 $q_{j}$ 的展开系数，$(\dd u/\dd a)_{j}$ 是周期导数（即低能
%  规范耦合的倒数）在奇点处的值。两个奇点相撞时，原来各自独立的局部展开被同一个“合并”
%  的局部坐标取代，$\kappa_{j}$ 与 $(\dd u/\dd a)_{j}$ 都随 $m\to m_{*}$ 趋于零（或获得
%  额外零点）。而 (11.3) 中二者的幂次分别是 $\chi_{h}$ 与 $\chi_{h}+\sigma$；当这些指数为负时，
%  零点就变成极点，这就是“可能出现发散”的来源。物理上不允许紧致流形上出现这种极点，
%  于是对 $\mathrm{SW}(\lambda)$ 施加了实质性的约束。]

% [Tutor Note: (11.4) 与小平分类的对照：原式只是在碰撞点附近的“正规形”，且只到首阶、
%  差数值系数。取 $A=z$、$B=\delta u$，判别式为 $\Delta=4A^{3}+27B^{2}=4z^{3}+27\delta u^{2}$，
%  其零点集是 $(z,\delta u)$ 平面中一条以原点为尖点的曲线——原本分离的两条奇点轨迹在此相切
%  并合并。若沿 $\delta u=0$ 切片，则 $\operatorname{ord}A=1$、$\operatorname{ord}B=\infty$、
%  $\operatorname{ord}\Delta=3$（小平 III 型）；若沿 $z=0$ 切片，则
%  $\operatorname{ord}A=\infty$、$\operatorname{ord}B=1$、$\operatorname{ord}\Delta=2$
%  （小平 II 型）。也就是说，两个 $I_{1}$ 型退化纤维碰撞之后不能再由单个 $I_{1}$ 描述，
%  而出现更退化的纤维型——这正是原文所说的 “further singularities”。由于 (11.4) 只是
%  首阶正规形，这些阶数结论只应在“差数值系数”的意义下理解。]

\subsection{超共形单型与广义诺特界}

对 $Z_{DW}(p,\Sigma;m)$ 的可能奇点作细致分析可以看出：当 $m\to m_{*}$ 时不出现发散，这一要求可由如下数学判据保证。

定义
\begin{equation}
  \mathrm{SW}_{X}(z):=\sum_{\lambda}
  \ee^{2\pi\ii\lambda_{c}\cdot\lambda}\,\mathrm{SW}(\lambda)\,\ee^{z\lambda},
  \tag{11.5}
\end{equation}
其中我们固定 $w_{2}(X)$ 的一个整提升 $2\lambda_{c}$，并把 $\lambda^{n}$ 的幂视为
$\mathrm{Sym}^{n}\bigl(\bar{H}^{2}(X)\bigr)$ 的对偶空间中的元素。于是

\begin{quote}
\itshape 若 $\mathrm{SW}_{X}(z)$ 在 $z=0$ 处解析，且具有不低于 $\chi_{h}-c-3$ 阶的零点，则
$Z_{DW}(p,\Sigma;m)$ 在 $m\to m_{*}$ 时有限。
\end{quote}

% [Tutor Note: 原文这条判据作 “with an of order $\ge\chi_{h}-c-3$”，其中 “an of order” 显然脱漏了
%  “zero” 一词（应为 “with a zero of order”）；“has a (nonnegative order) zero at $z=0$ of order
%  $\ge\cdots$” 一句也把 “nonnegative order” 夹在括号里，读起来别扭。译文按数学含义补足为
%  “具有不低于 $\chi_{h}-c-3$ 阶的零点”，并在此说明，以免读者误以为原文另有它意。]

我们\emph{定义} $X$ 为\emph{超共形单型}（superconformal simple type，SST），若 $\mathrm{SW}_{X}(z)$ 在 $z=0$ 处有一个（阶非负的）零点，其阶不低于 $\chi_{h}-c-3$。文献 \cite{ref33} 未能完全证明这也是 $Z_{DW}(p,\Sigma;m)$ 有限的必要条件，但已验证迄今一切可用的四维流形构造都满足这一判据。最近 \cite{ref18} 证明了射影簇都是 SST。

再往前推一点，就得到关于基本类数目的一个有趣下界。设 $B$ 是基本类的数目（其中两个非零类 $\lambda$ 与 $-\lambda$ 算作同一个）。则
\begin{equation}
  B\ge\Bigl\lfloor\frac{\chi_{h}-c}{2}\Bigr\rfloor,
  \tag{11.6}
\end{equation}
由此
\begin{equation}
  c\ge\chi_{h}-2B-1.
  \tag{11.7}
\end{equation}
代数几何的一个经典结果是：一般型极小曲面满足
\begin{equation}
  c\ge2\chi_{h}-6.
  \tag{11.8}
\end{equation}
这被称为诺特（Noether）界，因此我们把 (11.7) 称为“广义诺特界”。

这就为地理问题带来了一些新的界线：

\begin{figure}[H]
  \centering
  \figplaceholder{〔原文此处于第 71 页（PDF 第 72 页）留有图位；原 PDF 中未插入实际图形〕}
  \caption{$c$--$\chi_{h}$ 平面，其中画出了若干条适用 SW 求和规则的直线。}
  \label{fig:c-chih-plane}
\end{figure}
% [原 PDF 第 72 页]
% [Tutor Note: 原 PDF 此处也只有一行图题 “FIGURE: $c,\chi_{h}$ PLANE WITH SOME LINES WHERE
%  THE SW SUM RULES APPLY.”，并无图形。按上下文，图上应当画出 (11.7)（对不同的 $B$）与
%  (11.8) 这两族直线在 $c$--$\chi_{h}$ 平面中的位置，以显示 (11.7) 给出了诺特界之外的新界线；
%  本译本不作补画。]

% [Tutor Note: (11.5) 中的因子 $\ee^{2\pi\ii\lambda_{c}\cdot\lambda}$ 只是一个符号权重：因 $2\lambda_{c}$
%  是整类，$\lambda_{c}\cdot\lambda$ 取半整数值，故该因子恒为 $\pm1$。改变 $2\lambda_{c}$ 的整提升
%  等价于把 $\lambda$ 平移一个整类，此时 $\mathrm{SW}_{X}$ 只乘上 $\ee^{2\pi\ii\delta\cdot\lambda}$
%  型的相位（对偶环面上的平移），不改变其在 $z=0$ 附近的解析性质与零点阶，因此判据与提升的
%  选取无关。这正是原文说“固定一个整提升”却不必担心歧义的原因。]

% [Tutor Note: SST 判据的指数 $\chi_{h}-c-3$ 看似突兀，其实是 (11.3) 中各因子零点阶的净和：
%  在碰撞点附近，$\kappa_{j}^{\chi_{h}}$ 与 $(\dd u/\dd a)_{j}^{\chi_{h}+\sigma}$ 提供了负幂次
%  （当 $\chi_{h}<0$ 或 $\chi_{h}+\sigma<0$ 时），而 $\sum_{\lambda}\mathrm{SW}(\lambda)\ee^{z\lambda}$
%  在 $z=0$ 的零点阶恰是抵消这些负幂次所需的“裕量”。故判据的实质是：$\mathrm{SW}_{X}$ 在
%  $z=0$ 的零点阶必须不低于由 $\chi_{h}$、$\sigma$（即 $\chi$、$c$）决定的一个阈值。当
%  $\chi_{h}-c-3\le0$ 时条件自动满足，只需 $\mathrm{SW}_{X}$ 在 $0$ 处解析。原文在此只给出
%  充分性；必要性仍未证明（参见上引 [33] 的说明）。]

% [Tutor Note: 由 (11.6) 推 (11.7) 的关键是取整与整数性：$\lfloor x\rfloor\le B$ 等价于
%  $x<B+1$，即 $\chi_{h}-c<2B+2$，亦即 $c>\chi_{h}-2B-2$；因 $c$、$\chi_{h}$、$B$ 使右端为
%  整数（对 SW 单型，$c$ 与 $\chi_{h}$ 同奇偶性），故得 $c\ge\chi_{h}-2B-1$。这也反过来印证
%  (11.6) 用的是下取整 $\lfloor\cdot\rfloor$：若为上取整，相应的结论应是 $c\ge\chi_{h}-2B$，
%  与原文 (11.7) 不符。]

% [Tutor Note: 为什么把 $\lambda$ 与 $-\lambda$ 算作同一个“基本类”：共轭自旋$^{c}$ 结构给出
%  $\lambda\mapsto-\lambda$，而 (11.5) 的求和把两者分别计入，故 $B$ 是按 $\pm\lambda$ 配对后
%  的数目——这也解释了 (11.6) 中的除数 $2$。另可对比：若取 $B=1$（只有一个基本类对），
%  (11.7) 给出 $c\ge\chi_{h}-3$，当 $\chi_{h}>3$ 时弱于诺特界 (11.8) 的 $2\chi_{h}-6$；
%  (11.7) 的新意在于它对\emph{一切} SST 的光滑四维流形成立，而不限于代数曲面。
%  另：原文此处作 “(where we count two nonzero classes $\lambda$ and $-\lambda$ as the same.”，
%  括号未闭合（“Let $B$ the the number” 亦衍一 “the”），译文已补全、改正。]

\section{可能的未来方向}

\begin{keytakeaways}
\item \textbf{物理图像：} $u$ 平面积分并非终点：在 $b_{2}^{+}=0$ 的流形（如 $S^{3}\times S^{1}$）上，单圈项即已给出全部贡献，级数在一圈处截断；而把刚性超对称理论耦合到背景超引力、以及 class $\mathcal{S}$ 理论，则是把这一积分推广到更一般时空与理论的途径。
\item \textbf{数学工具：} 单圈项与 $u$ 平面上的截断、$S^{3}\times S^{1}$ 上 $\mathcal{N}=2$ 理论配分函数（Rastelli 等）、四维流形族的 $H^{*}(B\mathrm{Diff})$、class $\mathcal{S}$ 理论的交截理论。
\item \textbf{核心结论：} 作者留下四个方向与问题；最尖锐的是第 4 条：当 $\chi_{h}$ 为半整数时所有 SW 不变量为零，这质问物理是否对四维流形的“另一半”完全视而不见。
\end{keytakeaways}

\begin{enumerate}[label=\arabic*., leftmargin=2.1em]
\item 有一些有趣的情形，例如 $S^{3}\times S^{1}$，其中单圈项会有贡献。然而如 \cite{ref38} 所示，该级数在单圈处终止。最近，Rastelli 等人得到了 $S^{3}\times S^{1}$ 上 $\mathcal{N}=2$ 理论配分函数的漂亮结果。用 $u$ 平面积分重现这些结果将会很有意思。\footnote{〔原稿边注〕拓扑扭曲搞错了？}\footnote{〔原稿边注〕查一下。他们的指标中可能存在没有对应于唐纳森--威滕扭曲的特殊化。}

\item 四维流形的族与 $H^{*}(B\mathrm{Diff})$。最近关于把刚性超对称理论与背景超引力耦合的进展\footnote{〔原稿边注〕此处原稿只写了占位标记 “[CITE: SEIBERG et. al.]”。}应当会有所帮助。

\item 给出 class $\mathcal{S}$ 理论的“$u$ 平面积分”表达式。我们正在计算其交截理论的 UV 方程是什么？我们能否利用那一大批新的超共形理论学到新东西，也许沿着超共形单型这一思路？

\item 当 $\chi_{h}=\frac{1-b_{1}+b_{2}^{+}}{2}$ 为半整数时，所有 SW 不变量都为零。物理真的会对四维流形的半壁江山视而不见吗？
\end{enumerate}

% [Tutor Note: 第 1 条中“级数在单圈终止”是 $\S7$ 的一般结论在 $b_{2}^{+}=0$ 流形上的表现：
%  此时把度规标度 $t\to\infty$ 后，$u$ 平面上的瞬子修正被压掉，只剩单圈项（以及零模）。
%  但 $S^{3}\times S^{1}$ 的 $b_{1}=1$，故又须按 §10.3 的办法处理零模与通量扇区，
%  不能直接照搬 $b_{1}=0$ 的公式；这正是作者说“有趣”的原因。]

% [Tutor Note: 作者的两条边注（“WRONG TOPOLOGICAL TWIST?”与“他们的指标中可能存在没有对应于
%  唐纳森--威滕扭曲的特殊化”）针对的是同一件事：拓扑扭曲要求把 $SU(2)_{R}$ 的某个
%  $U(1)$ 子群与旋转群混合，而 Rastelli 等人所算的超共形指标通常是对另外的子群取的；
%  若两者所对应的指标没有公共特殊化，则不能指望用 $u$ 平面积分直接重现他们的结果。
%  这仍是原文未解决的问题，本译本照录边注而不代为判定。]

% [Tutor Note: 第 4 条是全篇最尖锐的物理诘问，原文只作断言而未给证明。其物理来源应当是
%  式 (9.1) 中的因子 $\ii^{\chi_{h}}$：$Z_{DW}$ 对实的 $(p,\Sigma)$ 应当是一个实/有理量
%  （它计数的是瞬子模空间的交截数），而 $\chi_{h}$ 为半整数时 $\ii^{\chi_{h}}$ 是非实的相位；
%  除非 SW 求和中的其它相位恰好将其抵消，否则只能有 $\mathrm{SW}(\lambda)\equiv0$。作者用
%  “物理是否对半个四维流形世界视而不见”来反问这一结论的合理性，暗示其中可能还有未被理解的
%  结构（例如需要更精细的 $b_{1}\neq0$ 处理，或某个离散 $\theta$ 角）。]
